% ==========================================================
%  preamble.tex --- Assume a Spherical Cow
% ==========================================================

% ---------------- fonts & encoding ----------------
\usepackage{amsmath,amsthm}
\usepackage{fontspec}
\usepackage[mathrm=sym]{unicode-math}
\setmainfont{texgyrepagella}[Extension=.otf,UprightFont=*-regular,BoldFont=*-bold,
  ItalicFont=*-italic,BoldItalicFont=*-bolditalic]
\setsansfont{texgyreheros}[Extension=.otf,UprightFont=*-regular,BoldFont=*-bold,
  ItalicFont=*-italic,BoldItalicFont=*-bolditalic,Scale=MatchLowercase]
\setmonofont{lmmono10-regular.otf}
\setmathfont{texgyrepagella-math.otf}
\usepackage{microtype}

% ---------------- page ----------------
\usepackage[letterpaper,top=1in,bottom=1in,inner=1.15in,outer=1.05in,headsep=18pt,headheight=15pt]{geometry}
\usepackage[dvipsnames]{xcolor}
\usepackage{enumitem}
\usepackage{booktabs}
\usepackage{tabularx}
\usepackage{graphicx}
\usepackage[most]{tcolorbox}
\usepackage{tikz}
\usetikzlibrary{arrows.meta,backgrounds,calc,positioning,decorations.pathmorphing,decorations.pathreplacing,decorations.markings,shapes.geometric}
\usepackage{pgfplots}
\pgfplotsset{compat=1.18,
  asc/.style={axis lines=left, axis line style={-{Stealth[length=5pt]}}, tick align=outside,
    tick label style={font=\footnotesize}, label style={font=\small},
    legend style={draw=none, fill=white, fill opacity=0.85, text opacity=1, font=\footnotesize},
    every axis plot/.append style={line join=round}}}
\usepackage{listings}
\usepackage{titlesec}
\usepackage{fancyhdr}

% ---------------- palette ----------------
\definecolor{ascblue}{HTML}{2B4EC0}
\definecolor{ascwarm}{HTML}{B23A2B}
\definecolor{ascpurple}{HTML}{7A3FB0}
\definecolor{ascbluebg}{HTML}{EEF1FC}
\definecolor{ascwarmbg}{HTML}{FAEEEB}
\definecolor{ascgraybg}{HTML}{F4F3EF}
\definecolor{ascrule}{HTML}{D9D6CC}
\definecolor{ascink}{HTML}{20242A}

\usepackage[colorlinks=true,linkcolor=ascblue,urlcolor=ascblue,citecolor=ascblue,bookmarksnumbered=true]{hyperref}

% ---------------- headings ----------------
\titleformat{\chapter}[display]
  {\normalfont\sffamily}
  {\large\color{ascblue}\MakeUppercase{\chaptertitlename}\ \thechapter}
  {4pt}
  {\Huge\bfseries\rmfamily\color{ascink}}
\titlespacing*{\chapter}{0pt}{10pt}{22pt}
\titleformat{\section}
  {\normalfont\Large\bfseries}{\color{ascblue}\S\thesection}{0.6em}{}
\titleformat{\subsection}
  {\normalfont\large\bfseries}{\color{ascblue}\thesubsection}{0.6em}{}
\setcounter{secnumdepth}{2}
\setcounter{tocdepth}{1}

% ---------------- running heads ----------------
\pagestyle{fancy}
\fancyhf{}
\renewcommand{\chaptermark}[1]{\markboth{#1}{}}
\renewcommand{\sectionmark}[1]{\markright{#1}}
\fancyhead[LE]{\small\sffamily\thepage\qquad\textit{Assume a Spherical Cow}}
\fancyhead[RO]{\small\sffamily\leftmark\qquad\thepage}
\renewcommand{\headrulewidth}{0pt}
\fancypagestyle{plain}{\fancyhf{}\fancyfoot[C]{\small\sffamily\thepage}\renewcommand{\headrulewidth}{0pt}}

% ---------------- theorem-like environments ----------------
\theoremstyle{definition}
\newtheorem{definition}{Definition}[chapter]
\newtheorem{example}[definition]{Example}
\newtheorem{exercise}[definition]{Exercise}
\newtheorem{question}[definition]{Question}
\newtheorem{theorem}[definition]{Theorem}
\newtheorem{law}[definition]{Law}
\newtheorem{proposition}[definition]{Proposition}
\theoremstyle{remark}
\newtheorem{remark}[definition]{Remark}

\tcolorboxenvironment{definition}{enhanced,breakable,colback=white,colframe=ascblue!55,
  boxrule=0.5pt,arc=2pt,left=7pt,right=7pt,top=5pt,bottom=5pt,before skip=10pt,after skip=10pt}
\tcolorboxenvironment{theorem}{enhanced,breakable,colback=ascbluebg,colframe=ascblue!55,
  boxrule=0.5pt,arc=2pt,left=7pt,right=7pt,top=5pt,bottom=5pt,before skip=10pt,after skip=10pt}
\tcolorboxenvironment{law}{enhanced,breakable,colback=ascbluebg,colframe=ascblue!55,
  boxrule=0.5pt,arc=2pt,left=7pt,right=7pt,top=5pt,bottom=5pt,before skip=10pt,after skip=10pt}
\tcolorboxenvironment{proposition}{enhanced,breakable,colback=ascbluebg,colframe=ascblue!55,
  boxrule=0.5pt,arc=2pt,left=7pt,right=7pt,top=5pt,bottom=5pt,before skip=10pt,after skip=10pt}
\tcolorboxenvironment{example}{enhanced,breakable,colback=ascgraybg,colframe=ascgraybg,
  boxrule=0pt,arc=2pt,left=7pt,right=7pt,top=5pt,bottom=5pt,before skip=10pt,after skip=10pt}

% prototype box (red) --- the "prototypical example" convention
\newtcolorbox{prototype}[1][PROTOTYPE]{enhanced,breakable,sharp corners,
  colback=ascwarmbg,colframe=ascwarmbg,boxrule=0pt,
  borderline west={3pt}{0pt}{ascwarm},
  left=12pt,right=10pt,top=7pt,bottom=7pt,before skip=10pt,after skip=12pt,
  fonttitle=\sffamily\bfseries\footnotesize,coltitle=ascwarm,
  attach title to upper={\quad},title={#1}}

% key-idea box (blue)
\newtcolorbox{keyidea}[1][KEY IDEA]{enhanced,breakable,arc=4pt,
  colback=ascbluebg,colframe=ascbluebg,boxrule=0pt,
  left=10pt,right=10pt,top=7pt,bottom=7pt,before skip=10pt,after skip=12pt,
  fonttitle=\sffamily\bfseries\footnotesize,coltitle=ascblue,
  attach title to upper={\quad},title={#1}}

% ---------------- chapter furniture ----------------
\newcommand{\buildson}[1]{\noindent{\small\sffamily\color{gray}{\footnotesize BUILDS ON}\enspace #1}\par\medskip}
\newcommand{\chapterdek}[1]{{\noindent\large\itshape #1\par}\bigskip}

% ---------------- difficulty: orders of perturbation ----------------
\definecolor{ordone}{HTML}{2B4EC0}
\definecolor{ordtwo}{HTML}{7A3FB0}
\definecolor{nonpert}{HTML}{B23A2B}
\newcommand{\diffmark}[1]{%
  \ifcase#1\or
    {\color{ordone}$\mathcal{O}(\epsilon)$}%
  \or
    {\color{ordtwo}$\mathcal{O}(\epsilon^{2})$}%
  \or
    {\color{nonpert}$e^{-1/\epsilon}$}%
  \fi}
\newenvironment{problems}{%
  \begin{enumerate}[label=\textbf{\thechapter.\arabic*},ref=\thechapter.\arabic*,
    leftmargin=3.1em,labelsep=0.6em,itemsep=9pt,topsep=6pt]}
  {\end{enumerate}}
\newcommand{\prob}[1]{\item \diffmark{#1}\quad}
\newcommand{\core}{\tikz[baseline=-0.55ex]\node[star,star points=5,star point ratio=2.3,
  fill=ascblue,inner sep=1.15pt]{};\ }
\newcommand{\aside}{\textcolor{gray}{$\dagger$}\ }
\newcommand{\hint}[1]{\par\smallskip{\small\textit{Hint:} #1}}

% ---------------- hints & solutions appendices ----------------
% \solchapter{qm} -> heading "14  Quantum mechanics"; \sol{qm}{5} -> "14.5"
\newcommand{\solchapter}[1]{\section*{{\color{ascblue}\ref*{ch:#1}}\quad\nameref*{ch:#1}}}
\newcommand{\sol}[2]{\par\smallskip\noindent\textbf{\ref*{ch:#1}.#2}\enspace}
\newcommand{\solfile}[1]{\IfFileExists{backmatter/solutions/#1.tex}{\input{backmatter/solutions/#1}}{}}

% ---------------- further reading ----------------
\newenvironment{reading}{%
  \section*{Further reading}
  \begin{itemize}[leftmargin=1.4em,itemsep=4pt]}
  {\end{itemize}}
\newcommand{\rd}[2]{\item \textbf{#1} --- #2}

% ---------------- code ----------------
\lstset{basicstyle=\ttfamily\small,columns=fullflexible,keepspaces=true,
  frame=single,rulecolor=\color{ascrule},backgroundcolor=\color{ascgraybg},
  xleftmargin=1em,framexleftmargin=0.6em,language=Python,
  keywordstyle=\color{ascblue}\bfseries,commentstyle=\color{gray}\itshape,
  stringstyle=\color{ascwarm},showstringspaces=false,breaklines=true}

% ---------------- math macros ----------------
\newcommand{\dd}{\mathop{}\!\mathrm{d}}
\newcommand{\dv}[2]{\frac{\mathrm{d}#1}{\mathrm{d}#2}}
\newcommand{\pdv}[2]{\frac{\partial #1}{\partial #2}}
\newcommand{\ket}[1]{\lvert #1\rangle}
\newcommand{\bra}[1]{\langle #1\rvert}
\newcommand{\braket}[2]{\langle #1\vert #2\rangle}
\newcommand{\expval}[1]{\langle #1\rangle}
\newcommand{\vb}[1]{\mathbf{#1}}
\newcommand{\vu}[1]{\hat{\mathbf{#1}}}
\newcommand{\un}[1]{\,\mathrm{#1}}
\newcommand{\order}[1]{\mathcal{O}\!\left(#1\right)}
\newcommand{\abs}[1]{\left\lvert #1\right\rvert}
\newcommand{\norm}[1]{\left\lVert #1\right\rVert}
\newcommand{\grad}{\nabla}
\newcommand{\curl}{\nabla\times}
\newcommand{\RR}{\mathbb{R}}
\newcommand{\CC}{\mathbb{C}}
\newcommand{\Msun}{M_{\odot}}
\newcommand{\Rsun}{R_{\odot}}
\newcommand{\Lsun}{L_{\odot}}
\newcommand{\kB}{k_{\mathrm{B}}}
\DeclareMathOperator{\tr}{tr}
\DeclareMathOperator{\sgn}{sgn}

% ---------------- misc ----------------
\setlength{\parskip}{0.35em}
\setlength{\parindent}{1.2em}
\raggedbottom
